\begin{frame}{Definition Gitter (Lattices)}
    \begin{definition}[Gitter]
        Sei $B = \{\mathbf{b}_1, \mathbf{b}_2, \dots, \mathbf{b}_m\}$ eine Menge linear unabhängiger Vektoren aus dem euklidischen Vektorraum $\mathbb{R}^n$. Dann heisst die Menge aller ganzzahligen Linearkombinationen
        \[
            \mathcal{L}(B) = \left\{ \sum_{i=1}^{m} a_i \mathbf{b}_i \;\middle|\; a_i \in \mathbb{Z} \right\} \subseteq \mathbb{R}^n
        \]
        das von $B$ erzeugte \textbf{Gitter} in $\mathbb{R}^n$. $B$ nennen wir die \textbf{Basis} des Gitters $\mathcal{L}$.
    \end{definition}
\end{frame}

\begin{frame}{Auffrischung: Euklidische Vektorräume}
    \begin{itemize}
        \item
    \end{itemize}
\end{frame}

\begin{frame}{Auffrischung: Lineare Unabhängigkeit}
    \begin{itemize}
        \item
    \end{itemize}
\end{frame}

\begin{frame}{Auffrischung: Linear Kombinationen}
    \begin{itemize}
        \item
    \end{itemize}
\end{frame}

\begin{frame}{Gitter als Matritze}
    \begin{itemize}
        \item Ein Gitter kann äquivalent aus einer Matrix $\mathbf{B}$ erzeugt werden, deren Spaltenvektoren die Basisvektoren sind:
        $\mathcal{L}(\mathbf{B}) = \{\mathbf{B}\mathbf{a} \mid \mathbf{a} \in \mathbb{Z}^m\}$
    \end{itemize}
\end{frame}
